% =============================================================================
% Appendix B --- Full Module Port Tables
% Derived verbatim from the port lists of parallel_interface.sv,
% controller.sv, and input_buffer.sv, and annotated with the
% "Implementing logic" column from Context 1 §6.
% =============================================================================

\chapter{Full Module Port Tables}
\label{app:port-tables}

This appendix tabulates every port of the three RTL modules with its
direction, width, and implementing logic (continuous assignment,
registered, or combinational with specific driver). Widths marked with
package constants are those the actual port declarations use.

\section{\modname{parallel\_interface}}
\label{app:ports-pi}

\begin{table}[htbp]
    \centering
    \caption{Full port list of \modname{parallel\_interface} (from
        \sv{source/Parallel\_Interface/parallel\_interface.sv}).}
    \label{tab:ports-pi}
    \begin{tabularx}{\linewidth}{l l c X}
        \toprule
        \textbf{Port} & \textbf{Dir} & \textbf{Width} & \textbf{Description / implementing logic} \\
        \midrule
        \signame{clk}       & in  & 1  & Clock; declared but \emph{not referenced} in the module body. \\
        \signame{reset}     & in  & 1  & Active-high reset; declared but \emph{not referenced} in the module body. \\
        \signame{host\_data} & in  & 32 & Host 32-bit packet: \sv{[31:24]} payload byte; \sv{[23:0]} one-hot
                                          tail \{PE, SA, col, row\}. \\
        \signame{host\_cmd}  & in  & 3  & Host command (\sv{cmd\_t} in \sv{parallel\_interface\_pkg}). \\
        \signame{valid}     & out & 1  & Legal-transaction qualifier; continuous assign
                                          \sv{(host\_cmd != CMD\_HIZ) \&\& ann\_tail\_is\_valid\_onehot(host\_data[23:0])}. \\
        \signame{data}      & out & 8  & Continuous assign \sv{host\_data[31:24]}. \\
        \signame{address}   & out & 16 & Continuous assign \sv{ann\_tail\_to\_parallel\_addr(host\_data[23:0])}. \\
        \signame{cmd}       & out & 3  & Continuous assign \sv{host\_cmd}. \\
        \bottomrule
    \end{tabularx}
\end{table}

\paragraph{Summary.}
All outputs are continuous assignments. \signame{valid} uses four
one-hot predicates ANDed together together with \signame{host\_cmd}
\(\neq\) \sv{CMD\_HIZ}.

\section{\modname{ann\_controller}}
\label{app:ports-controller}

The controller is parameterized by \sv{ADDR\_WIDTH} (default
\sv{controller\_pkg::DEFAULT\_ADDR\_WIDTH}~=~8), \sv{WEIGHT\_WIDTH}
(default \sv{DEFAULT\_WEIGHT\_WIDTH}~=~16), \sv{USE\_WEIGHT\_PULSE\_LUT}
(default \sv{1'b1}), and \sv{WEIGHT\_PULSE\_LUT\_FILE} (default
\sv{"target/Controller/programming\_inputs/weight\_pulse\_lut.mem"}).
Only the last two influence the generated logic.

\begin{table}[htbp]
    \centering
    \caption{Full port list of \modname{ann\_controller} (from
        \sv{source/Controller/controller.sv}).}
    \label{tab:ports-controller}
    \small
    \begin{tabularx}{\linewidth}{l l c X}
        \toprule
        \textbf{Port} & \textbf{Dir} & \textbf{Width} & \textbf{Description / implementing logic} \\
        \midrule
        \signame{clk}    & in  & 1 & Clock. \\
        \signame{rst\_n} & in  & 1 & Asynchronous active-low reset. \\
        \midrule
        \signame{valid}   & in & 1  & Accepted host-transaction qualifier from PI. \\
        \signame{data}    & in & 8  & Host payload byte from PI. \\
        \signame{address} & in & 16 & PI-decoded address \sv{\{6'b0, blk, sb, cid, rid\}}. \\
        \signame{cmd}     & in & \sv{CMD\_WIDTH} & Host command from PI (3 bits). \\
        \midrule
        \signame{ann\_reset}     & out & 1  & Asserted in \stateval{S\_RESET} combinational default. \\
        \signame{op\_done}       & in  & 1  & Core handshake, read in PROG and ERASE sub-FSMs as a
                                              level-sensitive combinational condition. \\
        \signame{ann\_core\_word} & out & 32 & Packed output; \sv{pack\_ann\_core\_word(...)} when
                                              \signame{state}\(\neq\)\stateval{S\_IDLE}, else all zeros. \\
        \signame{pulses}         & out & 3  & Pulse mode output: \sv{PULSE\_MODE\_*} or
                                              \sv{PULSE\_MODE\_HIZ} per combinational block. \\
        \signame{weight\_read\_data} & in & 4 & ADC/quantizer readback during verify. \\
        \midrule
        \signame{buf\_reg\_add}   & out & 6 & Buffer address (context-dependent). \\
        \signame{buf\_reg\_ctrl}  & out & 3 & Buffer control mode (\sv{CTRL\_IDLE},
                                              \sv{CTRL\_DATA\_LOAD}, \sv{CTRL\_WEIGHT\_READ},
                                              \sv{CTRL\_COMPUTE}, \sv{CTRL\_RESULT\_OUT}). \\
        \signame{buf\_read\_write} & out & 1 & 1\,=\,write, 0\,=\,read. \\
        \signame{buf\_bit\_sel}    & out & 3 & Bit index used by \signame{D0}--\signame{D7} during compute. \\
        \signame{buf\_data\_out}   & out & 8 & Captured byte for buffer writes. \\
        \signame{buf\_ready}       & in  & 1 & Buffer-side ready flag. \\
        \signame{buf\_data}        & in  & \sv{BUFFER\_DATA\_WIDTH} & Buffer byte-read port (8 bits). \\
        \midrule
        \signame{busy} & out & 1 & Status to higher-level; high in any non-idle state. \\
        \bottomrule
    \end{tabularx}
\end{table}

\paragraph{Implementing logic summary.}
Registered in \sv{always\_ff}: \signame{state}, \signame{prog\_state},
\signame{address\_reg}, \signame{data\_reg}, and counters
(\signame{pulse\_cnt}, \signame{data\_count}, \signame{row\_count},
\signame{bit\_count}, \signame{verify\_pulse\_idx}, \signame{retry\_cnt},
\signame{prog\_retry\_cnt}, \ldots). The combinational FSM block drives
\signame{next\_state}, buffer control, \signame{pulses}, and the
\signame{ann\_core\_word} packing path. \signame{weight\_read\_data} is
compared in \sv{VERIFY\_CHECK} and updates the retry registers on the
clock edge.

\section{\modname{input\_buffer}}
\label{app:ports-buffer}

\begin{table}[htbp]
    \centering
    \caption{Full port list of \modname{input\_buffer} (from
        \sv{source/Input\_Buffer/input\_buffer.sv}).}
    \label{tab:ports-buffer}
    \begin{tabularx}{\linewidth}{l l c X}
        \toprule
        \textbf{Port} & \textbf{Dir} & \textbf{Width} & \textbf{Description / implementing logic} \\
        \midrule
        \signame{clk}    & in  & 1 & Clock. \\
        \signame{rst\_n} & in  & 1 & Asynchronous active-low reset. \\
        \signame{data\_in} & in  & 8 & Write data from the parallel interface. \\
        \signame{ready}    & out & 1 & Mode/address-qualified ready; combinational. \\
        \signame{reg\_ctrl}    & in  & 3 & Buffer mode selector. \\
        \signame{buf\_read\_write} & in  & 1 & 1 = write, 0 = read. \\
        \signame{buf\_reg\_add}    & in  & 6 & Buffer address (base row for bit-serial lanes). \\
        \signame{bit\_sel}         & in  & 3 & Bit index (0..7) for lanes \signame{D0}--\signame{D7}. \\
        \signame{buf\_data}        & out & 8 & Byte-read; combinational mux over
                                               \sv{buffer\_reg[addr]}. \\
        \signame{D0} \ldots \signame{D7} & out & 1 & Bit-serial lanes;
                                               \signame{D}\(_k\)~=~\sv{buffer\_reg[addr+k][bit\_sel]} when
                                               \signame{read\_en} and \(\text{addr}{+}k\) is in range. \\
        \bottomrule
    \end{tabularx}
\end{table}

\paragraph{Implementing logic summary.}
\signame{write\_en} and \signame{read\_en} are continuous assigns;
\sv{buffer\_reg} updates in \sv{always\_ff} on
\signame{write\_en}~\sv{\&\&}~\sv{is\_valid\_addr(addr)}, and is cleared
to zero by the asynchronous reset loop. \signame{buf\_data},
\signame{D0}--\signame{D7}, and \signame{ready} are driven by
combinational \sv{\textasciigrave comb} blocks.
